\section{Second-order design}

With armature inductance \(L_a\), armature resistance \(R_a\), inertia \(J\),
viscous friction \(b\), and motor constant \(K_m\), the continuous-time
second-order transfer function is
\begin{equation}
  G(s)=
  \frac{K_m}
       {L_aJs^2+\left(L_ab+R_aJ\right)s+\left(R_ab+K_m^2\right)} .
  \label{eq:second-order-continuous-transfer}
\end{equation}
After zero-order-hold discretisation this is written as
\begin{equation}
  G(z)=\frac{b_1z+b_0}{z^2+a_1z+a_0}.
  \label{eq:second-order-discrete-transfer}
\end{equation}
The adaptive estimator therefore has four parameters, with regression vector
\begin{equation}
  \varphi(k)=
  \begin{bmatrix}
    -y(k-1) & -y(k-2) & u(k-1) & u(k-2)
  \end{bmatrix}.
  \label{eq:second-order-regressor}
\end{equation}

Example 7.6 of the reader does not use the free \(c(z)\) of the general
Diophantine system; it forces \(c(z)=z-1\), i.e. a PI controller. Our
second-order design does the same thing one order higher. Pure pole placement
fixes only the poles, not the static gain, and therefore leaves a steady-state
error; the simulated value was \(7.8\) percent. To remove this error, the
design forces
\begin{equation}
  c(z)=(z-1)(z+c_0),
  \label{eq:second-order-integrator}
\end{equation}
which gives \(T(1)=1\). Equation 7.16 of the reader gives the general design
equation whose linear system is solved for the controller coefficients at
every sampling instant.

The decisive point is the sampling time. For
\[
  K_m=0.05,\quad
  J=\SI{0.001}{kg.m^2},\quad
  b=0.0005,\quad
  L_a=\SI{2}{mH},\quad
  R_a=\SI{2}{\ohm},
\]
the continuous poles are \(-998.7\,\si{s^{-1}}\) for the electrical mode and
\(-1.75\,\si{s^{-1}}\) for the mechanical mode. The electrical time constant
is therefore about \(\SI{1}{ms}\). Table~\ref{tab:second-order-zoh-values}
shows the exact ZOH discretisation at the hardware sampling time and at a
sampling time that resolves the electrical pole.

\begin{scripttable}
  \centering
  \caption{Exact ZOH parameters of the second-order motor model.}
  \label{tab:second-order-zoh-values}
  \begin{tabular}{rrrrrrl}
    \toprule
    \(T\) & \(a_1\) & \(a_0\) & \(b_1\) & \(b_0\) & \multicolumn{2}{c}{\(z\)-poles} \\
    \midrule
    \(\SI{80}{ms}\) & \(-0.8692\) & \(0.000000\) & \(1.8467\) & \(0.0218\) & \(0.8692\) & \(0\) \\
    \(\SI{2}{ms}\)  & \(-1.1322\) & \(0.1352\)   & \(0.0284\) & \(0.0148\) & \(0.9965\) & \(0.1357\) \\
    \bottomrule
  \end{tabular}
\end{scripttable}

At \(T=\SI{80}{ms}\), the electrical pole is mapped to \(z=0\) and has fully
decayed between two samples. Two of the four parameters then no longer
contribute independent information, the regressor loses rank, and the plant is
structurally not identifiable. The closed loop can nevertheless still follow
the reference stably, even though the parameter estimate cannot converge under
that sampling, because the controller only needs certain combinations of the
parameters, not their individual values. A stable closed loop is therefore not
evidence of successful identification. At \(T=\SI{2}{ms}\), the poles are
\(0.9965\) and \(0.1357\), and all four parameters are identifiable.

The working design uses \(T=\SI{2}{ms}\), with PRBS dither added to the
manipulated variable. Clipping does not corrupt the estimate because the
regressor uses the post-saturation command. The estimator is used without
forgetting, \(\alpha=1\), because without persistent excitation \(\alpha<1\)
makes \(P\) grow without bound.

The desired closed-loop polynomial is
\begin{equation}
  q(z)=z(z-0.1357)(z-0.9608)^2 .
  \label{eq:second-order-desired-polynomial}
\end{equation}
The fast electrical pole is deliberately left in place in
\eqref{eq:second-order-desired-polynomial} because it is already much faster
than the desired bandwidth. If instead
\[
  q(z)=(z-0.9608)^4
\]
is demanded, the design slows the electrical pole by a factor of \(50\). The
controller must then invert the fast plant dynamics and becomes unstable
itself, with its pole at \(|z|=2.66\), and behind a saturation limit such a
controller necessarily diverges. The design rule is therefore to leave modes
that are much faster than the desired bandwidth where they are.
